%-------------------------------------------------------------------------------
% CAPITOLO 22
%-------------------------------------------------------------------------------
\chapter{Lezione 22}
\label{capitolo_22}
\textit{20/05/2026}

\section{Modello di Hopfield}
Ripartiamo dal modello di Hopfield.
Stiamo modellando l'attività di una rete neurale considerando un insieme di stati $\{S_i\}$, con $S_i = \pm 1$.
Indichiamo un pattern di memoria con $\{\xi_i^\mu\}$, dove $\xi_i^\mu = \pm 1$ per ogni neurone $i$, e $\mu$ è un indice che etichetta i diversi pattern memorizzati.
In analogia con il modello di Ising, possiamo definire l'Hamiltoniana:

\begin{equation} \label{eq:ch22_1}
H = -\frac{1}{2} \sum_{i, j} W_{ij}S_iS_j
\end{equation}

I pesi sinaptici $W_{ij}$ sono definiti come:

\begin{equation} \label{eq:ch22_2}
W_{ij} = W \sum_{\mu=1}^{K} \xi_i^\mu \xi_j^\mu
\end{equation}

Introducendo delle nuove variabili $S'_i = S_i \xi_i$, l'Hamiltoniana diventa:

\begin{equation} \label{eq:ch22_3}
H = -\frac{W}{2}\sum_{i,j} S'_i S'_j
\end{equation}

Da questa Hamiltoniana, possiamo derivare la probabilità di una configurazione di stati $\{S_i\}$ all'equilibrio termodinamico, data dalla distribuzione di Boltzmann:

\begin{equation} \label{eq:ch22_4}
P(\{S_i\}) = \frac{e^{-\beta H}}{Z}
\end{equation}

dove $Z$ è la funzione di partizione.
Se consideriamo un modello all'equilibrio descritto dalla statistica di Boltzmann, l'energia libera del sistema è data da:

\begin{equation} \label{eq:ch22_5}
F = -\frac{1}{\beta} \ln Z
\end{equation}

Nel limite $\beta \rightarrow \infty$ (bassa temperatura), il sistema tende a minimizzare l'energia $H$.
Il numero di pattern $K$ che possono essere immagazzinati e recuperati in modo affidabile deve essere inferiore al numero di neuroni $N$.
Si definisce il rapporto $\alpha = K/N$ come la "densità" dei pattern memorizzati. Si osserva che esiste un valore critico ($\alpha_c \approx 0.138$) al di sopra del quale la capacità di recupero crolla catastroficamente e la rete diventa incapace di distinguere i pattern memorizzati.

\vspace{0.5cm}
\noindent\fbox{%
    \parbox{\dimexpr\linewidth-2\fboxsep-2\fboxrule\relax}{%
        \textbf{!!}\\
        Un pattern di memoria è distribuito su tutta la rete. La memoria non è localizzata in un singolo neurone, ma emerge dall'interazione collettiva di tutti i neuroni. Questa caratteristica rende il sistema robusto: se un neurone collassa o viene distrutto, il comportamento generale della rete non cambia drasticamente.
    }%
}
\vspace{0.5cm}

Per quantificare quanto la configurazione attuale della rete $S = \{S_i\}$ assomigli a uno dei pattern memorizzati $\xi^\mu$, si introduce il concetto di overlap.

\begin{equation} \label{eq:ch22_6}
m_s = \max_{\mu} \left( \frac{1}{N} \sum_i \langle \xi_i^\mu s_i \rangle \right)
\end{equation}

Questa grandezza misura la propensione della rete a "ricordare" il pattern con cui ha la massima somiglianza all'equilibrio.
Vogliamo capire cosa succede se si distrugge un neurone.
Supponiamo di aver fissato un pattern $\mu = \bar{\mu}$ (quello per cui l'overlap è massimo) e che l'overlap sia vicino a 1:

\begin{equation} \label{eq:ch22_7}
m_s = \frac{1}{N} \sum_i \xi_i^{\bar{\mu}} S_i \sim 1
\end{equation}

Se il neurone $k$ viene distrutto, il nuovo overlap $m_s'$ diventa:

\begin{equation} \label{eq:ch22_8}
m_s' = \frac{1}{N} \sum_{i \neq k} \xi_i^{\bar{\mu}} S_i = m_s - \frac{1}{N} \xi_k^{\bar{\mu}} S_k
\end{equation}

La diminuzione è piccola, proporzionale a $1/N$.


\section{Dinamica di Apprendimento: Regola Hebbiana}
Le connessioni $W_{ij}$ possono evolvere nel tempo: $W_{ij} \equiv W_{ij}(t)$.
Se al tempo $t+1$ la configurazione della rete $\{S_i(t+1)\}$ corrisponde a un nuovo pattern $\xi_i^{K+1}$ da memorizzare, allora le connessioni si aggiornano.
Esiste una dinamica di apprendimento descritta dalla regola di Hebb:

\vspace{0.5cm}
\noindent\fbox{%
    \parbox{\dimexpr\linewidth-2\fboxsep-2\fboxrule\relax}{%
        \begin{equation} \label{eq:ch22_9}
        W_{ij} \rightarrow W_{ij} + W S_i(t+1) S_j(t+1)
        \end{equation}
    }%
}
\vspace{0.5cm}

Questa regola implica che la connessione tra due neuroni si rafforza se i loro stati sono correlati (entrambi attivi o entrambi inattivi). La forma base dei pesi $W_{ij}$ data precedentemente è essa stessa un risultato di una forma di apprendimento Hebbiano:

\begin{equation} \label{eq:ch22_10}
W_{ij} = W\sum_{\mu=1}^{K} \xi_i^\mu \xi_j^\mu
\end{equation}

Consideriamo due neuroni A e B. Supponiamo che le loro interazioni siano le stesse, il che non è vero in generale perché non c'è simmetria in una rete generale ($W_{BA} \neq W_{AB}$ a causa della diversa scelta di dendriti e assoni).
Grazie alla regola Hebbiana, possiamo schematizzare il comportamento della plasticità della memoria: l'interazione (efficacia sinaptica) tra neuroni che si attivano contemporaneamente tende ad aumentare progressivamente.


\section{Metodo della Massima Entropia}
Consideriamo un insieme di osservabili $\{O^\mu\}$ che siamo in grado di misurare sperimentalmente.
Sia $O_{EXP}^\mu$ il risultato delle nostre misure sperimentali.
Vogliamo capire qual è il modello minimale consistente con i valori sperimentali $\{O_{EXP}^\mu\}$.
Ogni possibile osservabile deve essere una funzione dello stato del sistema $S = \{S_i\}$, perciò:

\begin{equation*} 
O^\mu \equiv O^\mu(S) 
\end{equation*}

La distribuzione di probabilità dell'osservabile è tale che il valore medio teorico di ogni osservabile coincida con quello sperimentale:

\begin{equation} \label{eq:ch22_11}
\langle O^\mu(S) \rangle = O_{EXP}^\mu
\end{equation}

Il modello dovrebbe essere dedotto a partire da meno informazioni possibile, quindi dovrebbe essere quello a massima entropia. L'entropia associata alla distribuzione $P$ è:

\begin{equation} \label{eq:ch22_12}
S[P] = -\sum_S P(S) \ln P(S)
\end{equation}

\vspace{0.5cm}
\noindent\fbox{%
    \parbox{\dimexpr\linewidth-2\fboxsep-2\fboxrule\relax}{%
        Questa formulazione proviene, ad esempio, dall'ensemble microcanonico, dove la probabilità di uno stato $S$ è uniforme sugli stati con una data energia $E$:
        
        $P(S) = \frac{1}{\Gamma(E)} \quad \text{se } H(S) = E, \quad 0 \text{ altrimenti}$
        
        $\Gamma(E) = \sum_S \delta(H(S)-E)$ è il volume nello spazio delle fasi.
        L'entropia (con $k_B=1$) diventa:
        \begin{equation} \label{eq:ch22_13}
        S = -\sum_{S: H(S)=E} \frac{1}{\Gamma(E)} \ln\left(\frac{1}{\Gamma(E)}\right) = \ln \Gamma(E)
        \end{equation}
        
        Nell'ensemble canonico, invece, si ha:
        
        $P(S) = \frac{e^{-\beta H}}{Z}$ \quad con \quad $Z = e^{-\beta F}$
        
        L'entropia è:
        \begin{align} \label{eq:ch22_14}
        S[P] &= -\sum_S \frac{e^{-\beta H(S)}}{Z} [-\beta H(S) - \ln Z] \nonumber \\ 
        &= \beta (\langle E \rangle - F) \nonumber \\ 
        &= \frac{1}{T} (\langle E \rangle - F)
        \end{align}
        
        Questo è equivalente a dire che per una temperatura fissata si ha:
        \begin{equation} \label{eq:ch22_15}
        F = \langle E \rangle - TS
        \end{equation}
    }%
}
\vspace{0.5cm}

Torniamo alla massimizzazione dell'entropia $S[P] = -\sum_S P(S) \ln P(S)$.
Dobbiamo trovare $P(S)$. Massimizzando $S[P]$ senza vincoli si ha:

\begin{equation} \label{eq:ch22_16}
\frac{\delta S[P]}{\delta P(S)} = -\left[\ln P(S) + P(S) \cdot \frac{1}{P(S)}\right] = -[\ln P(S) + 1] = 0
\end{equation}

Questo implicherebbe:

\begin{equation} \label{eq:ch22_17}
\ln P(S) = -1 \rightarrow P(S) = \text{costante}  
\end{equation}

Da questo calcolo deduciamo che, se non sappiamo nulla del nostro sistema (=nessun vincolo sulle osservabili), il modello a massima entropia è quello in cui tutte le configurazioni sono equiprobabili, come nell'ensemble microcanonico.
Eseguiamo nuovamente il calcolo imponendo dei vincoli tramite i moltiplicatori di Lagrange.
Per il solo vincolo di normalizzazione $\sum_S P(S) = 1$:

\begin{equation} \label{eq:ch22_18}
\tilde{S}_{\lambda_0 }[P] = -\sum_S P(S) \ln P(S) - \lambda_0 \left(\sum_S P(S) - 1\right)
\end{equation}

Ponendo $\frac{\delta \tilde{S}}{\delta P(S)}= 0$ e $\frac{\partial \tilde{S}}{\partial \lambda_0} = 0$:

\begin{equation} \label{eq:ch22_19}
-[\ln P(S) + 1 + \lambda_0] = 0 \implies \ln P(S) = -1 - \lambda_0
\end{equation}

Da cui:

\begin{equation} \label{eq:ch22_20}
P(S) = \frac{1}{e^{1+\lambda_0}}
\end{equation}

Quindi:
\begin{itemize}
    \item Se abbiamo il minimo delle informazioni (=non imponiamo alcun vincolo), massimizzando l'entropia troviamo che c'è un'equiprobabilità costante.
    \item Se aggiungiamo un vincolo $\lambda_0$, otteniamo $P(S) = e^{-1-\lambda_0}$. Implementando la normalizzazione, $Z = e^{1+\lambda_0}$.
\end{itemize}

Se abbiamo $M$ vincoli dati dai valori medi sperimentali $\langle O^\mu(S) \rangle = O_{EXP}^\mu$, introduciamo $M$ moltiplicatori di Lagrange $\lambda_\mu$ oltre a $\lambda_0$ per la normalizzazione.

\begin{align} \label{eq:ch22_21}
\tilde{S}[P, \{\lambda_\mu\}] = &-\sum_S P(S) \ln P(S) - \lambda_0 \left(\sum_S P(S) - 1\right) \nonumber \\
&- \sum_{\mu=1}^M \lambda_\mu \left(\sum_S P(S) O^\mu(S) - O_{EXP}^\mu\right)
\end{align}

Derivando rispetto a $P(S)$ e ponendo a zero:

\begin{equation} \label{eq:ch22_22}
-[\ln P(S) + 1 + \lambda_0 + \sum_\mu \lambda_\mu O^\mu(S)] = 0
\end{equation}

Da cui si ottiene la distribuzione di probabilità a massima entropia:

\begin{equation} \label{eq:ch22_23}
P(S) = \frac{e^{-\sum_\mu \lambda_\mu O^\mu(S)}}{e^{1+\lambda_0}} = \frac{e^{-\sum_\mu \lambda_\mu O^\mu(S)}}{Z}
\end{equation}

Possiamo identificare il termine $\sum_\mu \lambda_\mu O^\mu(S)$ con un'Hamiltoniana generalizzata $H_{eff}(S)$.
Quindi:

\begin{equation} \label{eq:ch22_24}
P(S) = \frac{e^{-H_{eff}}}{Z}
\end{equation}

Questo modello risulta utile solo se si conoscono i parametri $\lambda_\mu$.
Chiaramente, il modello $P(S)$ che otteniamo dipende crucialmente dalla scelta delle osservabili $O^\mu(S)$ che misuriamo e includiamo come vincoli. È necessario scegliere il minor numero di parametri (e quindi di osservabili) sufficienti a descrivere il sistema.


\section{Applicazione Sperimentale}
Vediamo un esempio sperimentale: lo studio dell'attività neuronale in una porzione di cervello.
Si discretizza il tempo e il segnale elettrico di ciascun neurone $S_i(t) \in \{+1, -1\}$.
Sia $t_{obs}$ il tempo totale di osservazione. Per ogni istante di tempo discreto $t_k$, abbiamo una configurazione dell'attività di tutti i neuroni $S(t_k) = \{S_i(t_k)\}$.
L'obiettivo è determinare le frequenze di occorrenza delle diverse configurazioni $S$ e costruire un istogramma per stimare $P_{exp}(S)$. Tuttavia, estrarre la distribuzione di probabilità direttamente dai dati è molto problematico. Non funziona bene perché sono necessari tempi di osservazione estremamente lunghi per campionare adeguatamente gli eventi rari.

Le osservabili "buone", cioè quelle che possono essere stimate in modo affidabile da dati sperimentali limitati, sono quelle che coinvolgono momenti bassi della distribuzione. Ad esempio, la media $\langle x \rangle = \int dx P(x) x$ o il secondo momento $\langle x^2 \rangle = \int dx P(x) x^2$.
Quando si misurano queste osservabili sperimentalmente, è preferibile utilizzare stime dirette dei momenti piuttosto che tentare di ricostruire prima $P(x)$ e poi calcolare i momenti, poiché la stima di $P(x)$, specialmente nelle code, è difficile.
Quando l'ordine $m$ del momento $\langle x^m \rangle$ aumenta, l'approssimazione della sua stima peggiora a causa della crescente influenza delle code della distribuzione, che sono mal campionate.

Impostiamo il problema per i neuroni. Le osservabili più semplici che possiamo misurare sono i valori medi dell'attività dei singoli neuroni:

\begin{equation} \label{eq:ch22_25}
O^i(S) = \langle S_i \rangle \quad \text{per } i=1, \dots, N
\end{equation}

Il valore di aspettazione sperimentale per il neurone $i$ è:

\begin{equation} \label{eq:ch22_26}
\langle S_i \rangle_{exp} = \frac{1}{T_{obs}} \sum_{t=1}^{T_{obs}} S_i(t)
\end{equation}

Applicando il metodo della massima entropia con queste osservabili, la distribuzione di probabilità risultante è (assumendo $\beta=1$ e identificando $\lambda_i = -h_i$):

\begin{equation} \label{eq:ch22_27}
P(S) = \frac{1}{Z} e^{\sum_i h_i S_i}
\end{equation}

Questo è un modello di Ising in cui non ci sono interazioni tra coppie di neuroni.
I parametri $h_i$ sono determinati imponendo che i valori medi teorici coincidano con quelli sperimentali:

\begin{equation} \label{eq:ch22_28}
\langle S_i \rangle = \langle S_i \rangle_{exp}
\end{equation}

Per il modello di Ising senza interazioni, sappiamo che $\langle S_i \rangle = \tanh(h_i)$. Quindi:

\begin{equation} \label{eq:ch22_29}
\tanh(h_i) = \langle S_i \rangle_{exp}
\end{equation}

\begin{equation} \label{eq:ch22_30}
h_i = \text{arctanh}(\langle S_i \rangle_{exp})
\end{equation}

Quindi, direttamente dai dati misurati (le attività medie dei singoli neuroni), possiamo costruire i parametri $h_i$ del nostro modello a massima entropia.
Le misure del valor medio dell'attività di un singolo neurone $\langle S_i \rangle$ rappresentano la frequenza media di attività del neurone $i$.

\begin{equation} \label{eq:ch22_31}
\langle S_i \rangle = \frac{1}{T_{obs}} \sum_t S_i(t)
\end{equation}

Questa è una statistica a "un punto", cioè relativa a un singolo neurone, sommata su tutti i tempi.
Se si vuole andare oltre e fare una statistica a "due punti", si considerano due neuroni diversi, $i$ e $j$, e si calcola la loro correlazione:

\begin{equation} \label{eq:ch22_32}
C_{ij} = \langle S_i S_j \rangle
\end{equation}

Da questa si può ottenere la correlazione connessa:

\begin{equation} \label{eq:ch22_33}
C_{ij}^{conn} = \langle S_i S_j \rangle - \langle S_i \rangle \langle S_j \rangle
\end{equation}

Quindi, le osservabili che ci interessano sono, in prima approssimazione, il valor medio dell'attività dei singoli neuroni e, successivamente, le correlazioni tra coppie di neuroni.
Con queste informazioni (valori medi sperimentali $\langle S_i \rangle_{exp}$ e correlazioni sperimentali $\langle S_i S_j \rangle_{exp}$), vogliamo usare il metodo della massima entropia per derivare un modello che sia la più semplice distribuzione di probabilità che riproduca bene i momenti che conosciamo.

Se includiamo come vincoli sia le magnetizzazioni medie $m_i = \langle S_i \rangle_{exp}$ sia le correlazioni a coppie $C_{ij}^{exp} = \langle S_i S_j \rangle_{exp}$, il modello a massima entropia risultante avrà la forma:

\begin{equation} \label{eq:ch22_34}
P(S) = \frac{1}{Z} \exp\left( \sum_i h_i S_i + \sum_{i<j} J_{ij} S_i S_j \right)
\end{equation}

Questo è un modello di Ising con interazioni a coppie $J_{ij}$. I parametri $h_i$ e $J_{ij}$ sono i moltiplicatori di Lagrange associati ai vincoli su $\langle S_i \rangle$ e $\langle S_i S_j \rangle$ rispettivamente.
Se queste informazioni (a uno e due punti) non bastano per descrivere adeguatamente il sistema, si possono aggiungere ulteriori parametri, includendo informazioni a tre punti e così via. È cruciale che queste siano informazioni sperimentali.
